\section{組ごとの値と外れた組}\label{sup:values}
\begin{table}[tbp]
  \centering
  \small
  \caption{組ごとの値と組をまたいだ差}
  \label{tab:group_values}
  % paper/tables/group_values.tex -- generated by analyze/paper/tables_paper.py; sources in paper/tables/README.md
\begin{tabular}{@{}lll@{}}
\toprule
 & 自己側の積み上がり $g$ & 符号に従う分（緑好み） \\
\midrule
\multicolumn{3}{@{}l}{組ごとの値（最終5点の平均．経験なし / 経験50\% / 経験80\%）} \\
組1 & $+0.019$ / $+0.094$ / $+0.240$ & $-0.022$ / $+0.122$ / $+0.270$ \\
組2 & $+0.033$ / $+0.103$ / $+0.238$ & $-0.081$ / $+0.055$ / $+0.438$ \\
組3 & $+0.021$ / $+0.089$ / $+0.225$ & $-0.163$ / $+0.003$ / $+0.288$ \\
組4 & $+0.024$ / $+0.084$ / $+0.244$ & $-0.072$ / $+0.058$ / $+0.013$ \\
組5 & $+0.022$ / $+0.101$ / $+0.233$ & $+0.018$ / $-0.020$ / $+0.163$ \\
組6 & $+0.020$ / $+0.097$ / $+0.237$ & $-0.221$ / $+0.069$ / $+0.354$ \\
非共有（組1 / 組2 / 組3） & $+0.248$ / $+0.241$ / $+0.238$ & $-0.001$ / $-0.000$ / $-0.001$ \\
\multicolumn{3}{@{}l}{組をまたいだ差（6組の平均 [95\%信頼区間]，$t$ 分布，自由度5）} \\
経験50\% $-$ 経験なし & $+0.071$ [$+0.064$, $+0.079$] & $+0.138$ [$+0.028$, $+0.248$] \\
経験80\% $-$ 経験50\% & $+0.141$ [$+0.131$, $+0.152$] & $+0.207$ [$+0.050$, $+0.363$] \\
経験80\% $-$ 経験なし & $+0.213$ [$+0.205$, $+0.221$] & $+0.345$ [$+0.132$, $+0.557$] \\
\midrule[\heavyrulewidth]
 & 符号に従う分（好み混合） & 他者側の符号応答 G5b \\
\midrule
\multicolumn{3}{@{}l}{組ごとの値（最終5点の平均．経験なし / 経験50\% / 経験80\%）} \\
組1 & $-0.065$ / $+0.148$ / $+0.235$ & $-0.049$ / $+0.197$ / $+0.474$ \\
組2 & $-0.057$ / $+0.046$ / $+0.370$ & $-0.030$ / $+0.063$ / $+0.682$ \\
組3 & $-0.086$ / $-0.010$ / $+0.180$ & $-0.191$ / $+0.077$ / $+0.392$ \\
組4 & $-0.094$ / $+0.034$ / $+0.071$ & $+0.291$ / $+0.066$ / $+0.261$ \\
組5 & $+0.005$ / $+0.008$ / $+0.192$ & $+0.081$ / $+0.089$ / $+0.204$ \\
組6 & $-0.221$ / $+0.076$ / $+0.232$ & $+0.029$ / $+0.495$ / $+0.285$ \\
非共有（組1 / 組2 / 組3） & $-0.002$ / $+0.002$ / $-0.000$ & $-0.002$ / $-0.005$ / $-0.005$ \\
\multicolumn{3}{@{}l}{組をまたいだ差（6組の平均 [95\%信頼区間]，$t$ 分布，自由度5）} \\
経験50\% $-$ 経験なし & $+0.137$ [$+0.027$, $+0.246$] & $+0.143$ [$-0.109$, $+0.394$] \\
経験80\% $-$ 経験50\% & $+0.163$ [$+0.059$, $+0.267$] & $+0.218$ [$-0.066$, $+0.503$] \\
経験80\% $-$ 経験なし & $+0.300$ [$+0.174$, $+0.425$] & $+0.361$ [$+0.057$, $+0.665$] \\
\bottomrule
\end{tabular}

  % 表3：G3・符号に従う分（2データ）・G5b の，組1〜6 の経験なし・50%・80% の値（最終5点の平均），
  % 非共有条件（組1〜3）の値，組をまたいだ差（50%−経験なし，80%−50%，80%−経験なし）の平均と 95% CI。
  % 出典：v13 まとめ §3 表3・§6 表6・§11。
\end{table}

全6組では，経験80\%条件と経験なし条件の差は6組中5組で成り立ち，組をまたいだ差は $+0.361$ [$+0.057$, $+0.665$] で0を含まなかった．
一方，段階ごとの差は，経験50\%条件 $-$ 経験なし条件が $+0.143$ [$-0.109$, $+0.394$]，経験80\%条件 $-$ 経験50\%条件が $+0.218$ [$-0.066$, $+0.503$] で，どちらも0を含んだ．
したがって G5b で示せるのは，経験80\%条件の他者側の働きが経験なし条件より大きいことまでであり，段階的な強まりは示せていない．
外れた組は二つある．
組4では経験なし条件の G5b が $+0.291$ と大きく，経験80\%条件との差が出なかった．
組6では経験50\%条件（$+0.495$）が経験80\%条件（$+0.285$）を上回った（\ref{app:component}節）．
% 出典：v13 まとめ §10（(c) G5b 80%−経験なし 5/6），§11 表（G5b），§3 表3（s3・s5）。

組ごとに見ると，経験80\%条件と経験なし条件の差，および経験80\%条件と経験50\%条件の差は，どちらのデータでも6組中5組で成り立った（外れたのは組4）．
経験50\%条件と経験なし条件の差は，6組中3〜4組にとどまった．
% 出典：v13 まとめ §6 表6（符号に従う分の範囲，本番3組），§10 (b)(c)，§11 表。

ただし，この指標は組1〜4 では結果を見た後に定義したもので，事前登録してから計算したのは組5・6 だけである（本文3.4節）．
事前登録した関門 G5a（保った快・不快のもとでの他者側の傾きそのもの）は，経験80\%条件で本番の3組中1組しか合格しなかった．
G5a の傾きには，快・不快の符号に従う動きと，保つ前に予測の中に置いた A-2 の位置から本当の位置へ戻る動きが混ざっており，経験80\%条件では後者が前者を打ち消していたと読める（事後の分析による推測）．
符号に従う分は，この戻る動きを除くために定義した（\ref{app:prereg}節）．
% 出典：v13 まとめ §3 表1（G5a 1/3 組），§3.1（符号に依らない分が 80% で負 −0.38〜−0.05），執筆メモ §4 の H2b の注。

\subsection{他者側の指標で外れた組}
組4では，経験なし条件の G5b が $+0.291$ と大きく（ほかの組は $-0.191$〜$+0.081$），経験80\%条件（$+0.261$）との差が出なかった．
さらに組4の経験80\%条件は，符号に従う分が0前後であった（緑好みデータ $+0.013$，好み混合データ $+0.071$）．
快・不快は位置の情報も含むので，経験なし条件のモデルも快・不快を位置の手がかりとして学び，その働きが共有モジュールを通して process-2 にも及んだ可能性がある（推測）．
その強さが組によって違い，組4で強く出たと考えられる．

組6では，G5b は経験50\%条件（$+0.495$）の方が経験80\%条件（$+0.285$）より大きかったが，符号に従う分は経験80\%条件の方がずっと大きかった（緑好みデータ $+0.354$ 対 $+0.069$，好み混合データ $+0.232$ 対 $+0.076$）．
G5b は静止した A-2 の一つの範囲だけに人工的な入力を入れるのに対し，符号に従う分は A-2 の本当の快・不快の形を用いる．
組6の経験50\%条件のモデルは，一つの範囲への人工的な入力には強く反応するが，本当の快・不快の向きにはあまり従わないと読める（推測）．

組5では，G5b の経験80\%条件と経験50\%条件の差が5点中1点でしか成り立たなかった（経験80\%条件 $+0.204$，経験50\%条件 $+0.089$）．
% 出典：v13 まとめ §3 表3・注意（s3・s5），§10 表（s4 の G5b 80%−50% 1/5）。
